\documentclass[11pt]{article}
\usepackage[margin=1in]{geometry}
\usepackage{amsmath,amssymb,booktabs,graphicx,hyperref,natbib}
\usepackage[T1]{fontenc}
\hypersetup{colorlinks=true,linkcolor=blue,citecolor=blue,urlcolor=blue}
\newcommand{\NSeeds}{10}
\newcommand{\NTest}{2{,}650}
\newcommand{\FocusBps}{15}
\newcommand{\AvgPremium}{6.05}
\newcommand{\NFamily}{60}
\newcommand{\NRejectRW}{5}
\newcommand{\DiffAZeroFocus}{$-$0.14}
\newcommand{\CIAZeroFocus}{[$-$0.26, $-$0.01]}
\newcommand{\PRWAZeroFocus}{0.142}
\newcommand{\DiffAZeroFifty}{$-$0.31}
\newcommand{\CIAZeroFifty}{[$-$0.51, $-$0.09]}
\newcommand{\PRWAZeroFifty}{0.034}
\newcommand{\DiffAOneFocus}{$-$0.11}
\newcommand{\CIAOneFocus}{[$-$0.24, +0.03]}
\newcommand{\PRWAOneFocus}{0.370}
\newcommand{\DiffAOneFifty}{$-$0.30}
\newcommand{\CIAOneFifty}{[$-$0.52, $-$0.07]}
\newcommand{\PRWAOneFifty}{0.063}
\newcommand{\DiffAThreeFocus}{$-$0.08}
\newcommand{\CIAThreeFocus}{[$-$0.41, +0.23]}
\newcommand{\PRWAThreeFocus}{0.888}
\newcommand{\DiffAThreeFifty}{$-$0.24}
\newcommand{\CIAThreeFifty}{[$-$0.74, +0.24]}
\newcommand{\PRWAThreeFifty}{0.730}
\newcommand{\DiffAFourFocus}{$-$0.07}
\newcommand{\CIAFourFocus}{[$-$0.28, +0.12]}
\newcommand{\PRWAFourFocus}{0.837}
\newcommand{\DiffAFourFifty}{$-$0.25}
\newcommand{\CIAFourFifty}{[$-$0.64, +0.14]}
\newcommand{\PRWAFourFifty}{0.584}
\newcommand{\DiffAFiveFocus}{$-$0.14}
\newcommand{\CIAFiveFocus}{[$-$0.32, +0.03]}
\newcommand{\PRWAFiveFocus}{0.408}
\newcommand{\DiffAFiveFifty}{$-$0.37}
\newcommand{\CIAFiveFifty}{[$-$0.67, $-$0.07]}
\newcommand{\PRWAFiveFifty}{0.094}
\newcommand{\DiffLFocus}{$-$0.14}
\newcommand{\CILFocus}{[$-$0.23, $-$0.04]}
\newcommand{\PRWLFocus}{0.033}
\newcommand{\DiffLFifty}{$-$0.18}
\newcommand{\CILFifty}{[$-$0.26, $-$0.07]}
\newcommand{\PRWLFifty}{0.004}
\newcommand{\DiffHWFocus}{+0.15}
\newcommand{\CIHWFocus}{[+0.06, +0.26]}
\newcommand{\PRWHWFocus}{1.000}
\newcommand{\DiffHWFifty}{+0.17}
\newcommand{\CIHWFifty}{[+0.07, +0.28]}
\newcommand{\PRWHWFifty}{1.000}
\newcommand{\DiffATwobFocus}{+0.10}
\newcommand{\CIATwobFocus}{[$-$0.14, +0.33]}
\newcommand{\PRWATwobFocus}{0.997}
\newcommand{\DiffATwobFifty}{$-$0.06}
\newcommand{\CIATwobFifty}{[$-$0.47, +0.34]}
\newcommand{\PRWATwobFifty}{0.931}
\newcommand{\PRWLFocusTwoM}{0.017}
\newcommand{\PRWLFocusQ}{0.064}


\title{Does Deep Hedging Survive an Honest Evaluation?\\
\large Training Data, Evaluation Protocol and Transaction Costs on Hong Kong Technology Stocks}
\author{Apiwit Triwatana\thanks{Draft for discussion with a faculty supervisor. Code and data:
\url{https://github.com/apiwittri2007/DeepHedging_Project}.}}
\date{Working draft, November 2026}

\begin{document}
\maketitle

\begin{abstract}
Deep hedging trains a neural network to hedge a derivative by directly minimising a risk measure
of hedging P\&L, including trading costs. Most studies train it on simulated paths and evaluate it
on a single test split. We ask two questions on 29 Hang Seng Tech stocks (2000--2026):
which training data makes a better deep hedger, and does the answer survive an honest evaluation?
We compare six training-data sources (GBM with known or estimated volatility, GARCH-$t$, a
stationary block bootstrap, real windows, and GBM pre-training fine-tuned on real windows) with
Black--Scholes delta and two three-parameter linear corrections to delta (a quadratic-in-delta
regression in the spirit of Ruf and Wang, and the Hull--White minimum-variance form). Every model is retrained each year using
only past data (walk-forward, 2016--2026, \NTest{} one-month option windows), with \NSeeds{} training
seeds, month-clustered bootstrap inference and a Romano--Wolf correction over the
\NFamily{} arm-by-cost comparisons.
Three findings stand out. First, a gain over delta that is significant in a single train/test split with one
seed is not robust under walk-forward retraining with several seeds. Second, at realistic Hong Kong
costs (\FocusBps{} bps per side) the quadratic-in-delta correction is as good as the best deep hedger and
the only method significant after correction, although only borderline when the bootstrap clusters
quarters instead of months.
Third, deep hedgers only pull ahead when costs are high (50 bps). More realistic training data does not
help beyond noise. The evaluation protocol matters more than the training data.
\end{abstract}

\section{Introduction}
Black--Scholes (BS) delta hedging is optimal only with continuous, costless trading and a correct
volatility. Deep hedging \citep{buehler2019deep} drops these assumptions: a neural network maps
market features to hedge positions and is trained to minimise a convex risk measure such as
CVaR of the terminal hedging P\&L, with transaction costs inside the objective. Since the network
needs many paths, it is usually trained on a simulator; how good that simulator must be, and whether real
data should be used instead, is an open question
\citep{mikkila2023empirical,buehler2022drift,hashimoto2026realism}.

Two strands of evidence motivate this paper. \citet{ruf2022hedging} show that once the evaluation
is done carefully, a linear regression on Black--Scholes sensitivities (delta, vega, vanna, gamma)
hedges as well as neural networks, and
\citet{ruf2020review} warn that conclusions often depend on the data split and on the volatility estimate.
Recent work finds that deep hedgers learn a delta reduction whose direction is consistent with the
minimum-variance delta of \citet{hull2017optimal} and that is fragile across regimes
\citep{zernikov2026learn}, that their advantage over simpler hedges grows with transaction costs
\citep{francois2025ivs,bracha2025drl}, and that the difference between deep and delta hedging can be
read as a statistical arbitrage \citep{francois2024statarb}. Architectures and features have been
extended in many directions, from recurrent networks on real index and Hong Kong options
\citep{zhang2021option} to no-transaction-band networks \citep{imaki2021ntb}, reinforcement learning
\citep{cao2021deep} and residual corrections to delta \citep{qiao2024residual}; generative models such
as \citet{wiese2020quant} supply realistic training paths. Closest to our training-data question,
\citet{mikkila2023empirical} find that training on real S\&P 500 option prices beats training on
calibrated simulations, whereas our real-window arm does not beat GBM training; with synthetic option
prices (Section~\ref{sec:limits}) the two results are not directly comparable.

We contribute in three ways.
(i) \emph{Protocol fragility.} The same deep hedger that beats delta significantly on one split with
one seed loses that significance under yearly walk-forward retraining with \NSeeds{} seeds; the
seed-to-seed spread is as large as the effect.
(ii) \emph{Training-data comparison under one protocol.} Six training-data sources, pre-registered before
results were seen, are evaluated on the same \NTest{} real test windows; realism of the training data
(tails, volatility clustering) does not translate into better hedging.
(iii) \emph{A high-cost emerging market.} Hong Kong single stocks carry stamp duty of 10 bps per side;
we map where deep hedging pays off (high costs) and where a three-parameter correction is enough
(realistic costs).

\section{Data}
Daily adjusted closes for 29 Hang Seng Tech constituents (2000--2026) come from a public GitHub
mirror of Yahoo Finance data. We study one-month (21 trading days) at-the-money European calls,
priced and hedged with the 63-day trailing realised volatility known at inception (no look-ahead).
Windows with a daily log return above 50\% in absolute value are treated as data errors.
The test set consists of non-overlapping windows starting in 2016--2026 (\NTest{} windows).
\emph{Limitation:} option prices are synthetic (BS with trailing volatility), not traded quotes;
Section~\ref{sec:limits} discusses this.

\section{Method}
\paragraph{Hedging problem.} A trader sells one call (strike $K=100$, spot $S_0=100$, $T=21/252$)
for its BS price and rebalances the underlying daily, paying a proportional cost $c$ per unit traded.
P\&L is the self-financing portfolio value minus the payoff. We report CVaR$_{95}$ of P\&L
(average loss in the worst 5\% of windows) in HK\$ per HK\$100 of underlying.

\paragraph{Deep hedger.} A 4--32--32--1 MLP maps (time to maturity, log-moneyness, current holding,
BS delta) to the new holding, and is trained to minimise CVaR$_{95}$ in the Rockafellar--Uryasev form
\citep{rockafellar2000optimization}, differentiating through the whole path.

\paragraph{Training-data arms.} For each test year $Y$ every arm uses only prices dated before
1 January $Y$; real training windows must also end before $Y$.
A0: GBM with the trader knowing the true volatility; A1: GBM where the trader estimates volatility
from a 63-day warm-up, exactly as on real data; A2: GARCH(1,1)-$t$ fitted per stock (A2b, added after
the main run and labelled post hoc: Student-$t$ degrees of freedom $\nu\ge 4$); A3: stationary block
bootstrap of real returns \citep{politis1994stationary}; A4: real overlapping windows; A5: A1 fine-tuned on A4.

\paragraph{Benchmarks.} BS delta; delta with a no-trade band; no hedge; L, the linear correction
$h=\delta+a+b\delta+c\delta^2$; and HW, the \citet{hull2017optimal} minimum-variance form
$h=\delta+n(d_1)(a+b\delta+c\delta^2)$. Without traded option prices, L and HW coefficients are
fitted by least squares to the variance of zero-cost hedging P\&L on the same pre-$Y$ real windows.

\paragraph{Inference.} CVaR is averaged over \NSeeds{} seeds. Differences with delta are tested with a
paired bootstrap that resamples whole calendar months (2{,}000 draws), so stocks hit by the same
shock are not counted as independent. All \NFamily{} arm-by-cost comparisons form one family:
we report Holm \citep{holm1979simple} and Romano--Wolf step-down \citep{romano2005stepwise}
adjusted $p$-values (statistics standardised by each test's bootstrap standard error) from a single
shared bootstrap \citep[see also][]{harvey2016and}; $p$-values use $(1+\#)/(B+1)$. Because a 21-day window
can span two calendar months, we repeat the inference with 2- and 3-month clusters (Table~4).
Inference treats the \NSeeds{} trained networks as given: it is conditional on the seeds.

\section{Results}
\begin{table}[h]\centering\small
\caption{CVaR$_{95}$ of hedging P\&L on the real test set (mean over \NSeeds{} seeds; lower is better).}
\begin{tabular}{lrrrrrr}
\toprule
Arm & 0 bps & 5 bps & 10 bps & 15 bps & 25 bps & 50 bps \\
\midrule
A0 GBM, known vol & 5.82 & 5.95 & 6.12 & 6.21 & 6.50 & 7.16 \\
A1 GBM, estimated vol & 5.86 & 5.96 & 6.07 & 6.24 & 6.51 & 7.16 \\
A2 GARCH-$t$ & 6.53 & 10.26 & 7.83 & 7.51 & 6.98 & 10.58 \\
A2b GARCH-$t$, $\nu\ge4$$^\dagger$ & 6.00 & 6.16 & 6.23 & 6.45 & 6.77 & 7.41 \\
A3 Block bootstrap & 5.84 & 5.97 & 6.12 & 6.28 & 6.52 & 7.23 \\
A4 Real windows & 5.90 & 6.02 & 6.17 & 6.29 & 6.58 & 7.22 \\
A5 GBM $\to$ real & 5.82 & 5.95 & 6.07 & 6.21 & 6.47 & 7.10 \\
L Linear correction & 5.77 & 5.92 & 6.06 & 6.21 & 6.52 & 7.29 \\
HW Min.-variance delta & 6.03 & 6.19 & 6.35 & 6.50 & 6.82 & 7.63 \\
BS delta & 5.89 & 6.04 & 6.20 & 6.35 & 6.67 & 7.47 \\
Delta + band & 5.91 & 6.06 & 6.21 & 6.36 & 6.66 & 7.43 \\
No hedge & 34.43 & 34.43 & 34.43 & 34.43 & 34.43 & 34.43 \\
\bottomrule
\end{tabular}

\end{table}

\begin{table}[h]\centering\small
\caption{CVaR difference vs BS delta at \FocusBps{} and 50 bps. Month-clustered 95\% CI; one-sided $p$-values,
raw and adjusted over all \NFamily{} comparisons. $^\dagger$post hoc.}
\begin{tabular}{lrrlrrr}
\toprule
Arm & Cost & $\Delta$CVaR & 95\% CI & $p$ & $p_{\text{Holm}}$ & $p_{\text{RW}}$ \\
\midrule
A0 GBM, known vol & 15 & $-$0.143 & [$-$0.26, $-$0.01] & 0.018 & 0.943 & 0.142 \\
A1 GBM, estimated vol & 15 & $-$0.112 & [$-$0.24, +0.03] & 0.061 & 1.000 & 0.370 \\
A2 GARCH-$t$ & 15 & +1.158 & [+0.51, +1.76] & 1.000 & 1.000 & 1.000 \\
A2b GARCH-$t$, $\nu\ge4$$^\dagger$ & 15 & +0.100 & [$-$0.14, +0.33] & 0.819 & 1.000 & 0.997 \\
A3 Block bootstrap & 15 & $-$0.076 & [$-$0.41, +0.23] & 0.326 & 1.000 & 0.888 \\
A4 Real windows & 15 & $-$0.067 & [$-$0.28, +0.12] & 0.251 & 1.000 & 0.837 \\
A5 GBM $\to$ real & 15 & $-$0.139 & [$-$0.32, +0.03] & 0.057 & 1.000 & 0.408 \\
HW Min.-variance delta & 15 & +0.152 & [+0.06, +0.26] & 0.998 & 1.000 & 1.000 \\
L Linear correction & 15 & $-$0.139 & [$-$0.23, $-$0.04] & 0.004 & 0.228 & 0.033 \\
Delta + band & 15 & +0.005 & [$-$0.08, +0.09] & 0.541 & 1.000 & 0.977 \\
A0 GBM, known vol & 50 & $-$0.306 & [$-$0.51, $-$0.09] & 0.003 & 0.174 & 0.034 \\
A1 GBM, estimated vol & 50 & $-$0.304 & [$-$0.52, $-$0.07] & 0.006 & 0.336 & 0.063 \\
A2 GARCH-$t$ & 50 & +3.115 & [+2.41, +4.23] & 1.000 & 1.000 & 1.000 \\
A2b GARCH-$t$, $\nu\ge4$$^\dagger$ & 50 & $-$0.060 & [$-$0.47, +0.34] & 0.402 & 1.000 & 0.931 \\
A3 Block bootstrap & 50 & $-$0.240 & [$-$0.74, +0.24] & 0.181 & 1.000 & 0.730 \\
A4 Real windows & 50 & $-$0.246 & [$-$0.64, +0.14] & 0.106 & 1.000 & 0.584 \\
A5 GBM $\to$ real & 50 & $-$0.367 & [$-$0.67, $-$0.07] & 0.011 & 0.583 & 0.094 \\
HW Min.-variance delta & 50 & +0.167 & [+0.07, +0.28] & 1.000 & 1.000 & 1.000 \\
L Linear correction & 50 & $-$0.179 & [$-$0.26, $-$0.07] & 0.001 & 0.060 & 0.004 \\
Delta + band & 50 & $-$0.040 & [$-$0.12, +0.05] & 0.226 & 1.000 & 0.746 \\
\bottomrule
\end{tabular}

\end{table}

\begin{table}[h]\centering\small
\caption{Romano--Wolf adjusted $p$-values under 1-, 2- and 3-month bootstrap clusters
(rows with any $p<0.10$). The pre-registered choice is 1 month.}
\begin{tabular}{lrrrrr}
\toprule
Arm & Cost & $\Delta$CVaR & $p_{\text{RW}}$ 1-month & 2-month & 3-month \\
\midrule
L Linear correction & 0 & $-$0.118 & 0.094 & 0.048 & 0.170 \\
L Linear correction & 5 & $-$0.126 & 0.068 & 0.034 & 0.125 \\
L Linear correction & 10 & $-$0.132 & 0.049 & 0.024 & 0.091 \\
L Linear correction & 15 & $-$0.139 & 0.033 & 0.017 & 0.064 \\
L Linear correction & 25 & $-$0.150 & 0.018 & 0.010 & 0.042 \\
A0 GBM, known vol & 50 & $-$0.306 & 0.034 & 0.053 & 0.042 \\
A1 GBM, estimated vol & 50 & $-$0.304 & 0.063 & 0.061 & 0.049 \\
A5 GBM $\to$ real & 50 & $-$0.367 & 0.094 & 0.076 & 0.109 \\
L Linear correction & 50 & $-$0.179 & 0.004 & 0.004 & 0.017 \\
\bottomrule
\end{tabular}

\end{table}

\begin{table}[h]\centering\small
\caption{Economic size vs BS delta (average premium HK\$\AvgPremium{} per HK\$100).}
\resizebox{\textwidth}{!}{\begin{tabular}{lrrrr}
\toprule
Arm & Cost & $\Delta$CVaR (HK\$ per HK\$1m) & \% of premium & $\Delta$mean P\&L (HK\$ per HK\$1m) \\
\midrule
A0 GBM{,} known vol & 15 & $-$1{,}428 & $-$2.4 & $-$67 \\
A1 GBM{,} estimated vol & 15 & $-$1{,}123 & $-$1.9 & $-$14 \\
A2 GARCH-$t$ & 15 & +11{,}582 & +19.2 & +46 \\
A2b GARCH-$t${,} $\nu\ge4$$^\dagger$ & 15 & +1{,}004 & +1.7 & +91 \\
A3 Block bootstrap & 15 & $-$763 & $-$1.3 & +282 \\
A4 Real windows & 15 & $-$670 & $-$1.1 & +369 \\
A5 GBM $\to$ real & 15 & $-$1{,}386 & $-$2.3 & +469 \\
HW Min.-variance delta & 15 & +1{,}522 & +2.5 & +138 \\
L Linear correction & 15 & $-$1{,}387 & $-$2.3 & +240 \\
Delta + band & 15 & +52 & +0.1 & +317 \\
A0 GBM{,} known vol & 50 & $-$3{,}057 & $-$5.1 & +1{,}066 \\
A1 GBM{,} estimated vol & 50 & $-$3{,}040 & $-$5.0 & +1{,}264 \\
A2 GARCH-$t$ & 50 & +31{,}148 & +51.5 & $-$1{,}182 \\
A2b GARCH-$t${,} $\nu\ge4$$^\dagger$ & 50 & $-$601 & $-$1.0 & +1{,}569 \\
A3 Block bootstrap & 50 & $-$2{,}399 & $-$4.0 & +1{,}594 \\
A4 Real windows & 50 & $-$2{,}456 & $-$4.1 & +1{,}649 \\
A5 GBM $\to$ real & 50 & $-$3{,}668 & $-$6.1 & +1{,}696 \\
HW Min.-variance delta & 50 & +1{,}671 & +2.8 & $-$64 \\
L Linear correction & 50 & $-$1{,}793 & $-$3.0 & +254 \\
Delta + band & 50 & $-$396 & $-$0.7 & +917 \\
\bottomrule
\end{tabular}
}
\end{table}

% Results prose: paper/sections/results.tex (numbers come from generated/numbers.tex)
\paragraph{Protocol fragility.} The pilot study trained one deep hedger (arm A0) on a single
pre-2016 split with one seed. It reported a CVaR reduction of $-0.25$ at \FocusBps{} bps, significant
under the month-clustered bootstrap. Under walk-forward retraining the same arm gives \DiffAZeroFocus{}
(95\% CI \CIAZeroFocus{}) with \NSeeds{} seeds, and only $-0.05$ with the first three seeds.
Its seed-to-seed standard deviation is 0.05--0.15 across arms, the same order as the effect.
After the Romano--Wolf correction over all \NFamily{} comparisons its adjusted $p$-value is \PRWAZeroFocus{}.
A single split with a single seed therefore overstates both the size and the reliability of the gain.

\paragraph{Realistic costs: a three-parameter correction is enough.} At \FocusBps{} bps the linear
correction L reduces CVaR by \DiffLFocus{} (CI \CILFocus{}, Romano--Wolf $p=$\PRWLFocus{}). It beats delta
in 10 of 11 test years and is the only method significant after correction at 10, 15 and 25 bps.
The best deep arms (A0 and A5) reach the same point estimate, but with wider intervals and no
significance after correction (Table~2). In money terms L saves about HK\$1{,}400 of tail loss per HK\$1m of
option notional, 2.3\% of the average premium (Table~3).

\paragraph{High costs: deep hedging pays.} At 50 bps the deep arms trained on GBM
(\DiffAZeroFifty{} for A0, \DiffAOneFifty{} for A1) and A5 (\DiffAFiveFifty{}) cut tail risk by 0.30--0.37,
about 5--6\% of the premium, roughly twice as much as L (\DiffLFifty{}). A0 remains significant after the
Romano--Wolf correction ($p=$\PRWAZeroFifty{}); A1 and A5 are borderline (\PRWAOneFifty{}, \PRWAFiveFifty{}).
This matches reports that deep-hedging gains grow with transaction costs \citep{francois2025ivs,bracha2025drl}.

\paragraph{Training data: realism does not help.} The block bootstrap reproduces the real data's tails
and volatility clustering closely (excess kurtosis 7.7 vs 6.7; Hill index 3.4 vs 3.5), while GBM has
none. Yet training on it (A3: \DiffAThreeFocus{}) or on real windows (A4: \DiffAFourFocus{}) is no better
than training on GBM. GARCH-$t$ fitted to pre-year data overshoots the tails (kurtosis 136) and makes
training unstable; capping the tails (A2b, post hoc) stabilises it but does not beat delta. This agrees
with \citet{hashimoto2026realism}: fidelity of the generator is not what drives hedging performance.

\paragraph{Hull--White form.} Fitted to the variance of hedging P\&L, the minimum-variance form
of \citet{hull2017optimal} \emph{increases} CVaR relative to delta at every cost. Its $n(d_1)$ scaling
puts the correction where the option's gamma is largest. Without traded option prices, there is no
spot--implied-volatility co-movement for it to capture, so the variance-optimal fit does not carry over
to the tail. We read this as a limitation of synthetic option prices rather than evidence against Hull--White.


\section{Limitations and next steps}\label{sec:limits}
Option prices are synthetic, so the spot--implied-volatility co-movement that drives the empirical
minimum-variance correction \citep{hull2017optimal,zernikov2026learn} is absent; a replication on
traded options (Deribit BTC options, and HKEX stock options if access is obtained) is in progress.
Constituents are today's index members (survivorship bias), and prices are dividend-adjusted, so the
option is effectively written on a total-return series. Test windows with a missing price or a daily
move above 50\% are dropped, which uses information from inside the window (one test window). We hedge single stocks with the stock
itself; index futures would be cheaper. Hyper-parameters were fixed before the walk-forward results
were seen and not tuned.

\section{Conclusion}
On \NTest{} real one-month option windows on Hong Kong technology stocks, an honest protocol changes
the conclusion. That protocol has three parts: yearly walk-forward retraining, \NSeeds{} seeds, and
inference that is month-clustered and corrected for multiple testing. A deep hedger that looks clearly
better on one split is not reliably better at realistic costs, where a three-parameter linear correction to delta
is as good and more robust. Deep hedging earns its keep only when trading is expensive. More
realistic training data does not help. For practitioners in high-cost markets, the evaluation protocol
and the cost level matter more than the training data. Replicating these findings on traded option
prices is the natural next step.


\bibliographystyle{plainnat}
\bibliography{references}
\end{document}
